% Comando simples para exibir comandos Latex no texto
\newcommand{\comando}[1]{\textbf{$\backslash$#1}}

\section{Considerações Gerais}

Testando a inserção de uma nova sigla \sigla{EESC}{Escola de Engenharia de São Carlos}, vamos ver se deu certo.

\sigla{ESR}{Estrutura de Solo Reforçado}

\sigla{RCD}{Resíduo de Construção e Demolição}

Testando citação \cite{NBR6457:1986}


\begin{table}[htb]
\IBGEtab{%
  \caption{População dos países da América do Sul} \label{tabela:populacao_america_sul_nova}
}{%
\begin{tabular}{r|p{3cm}|r}        
\toprule
Código  & País            & População   \\ \midrule \midrule
1       & Brasil          & ~~~~~~191.480.630 \\ \midrule 
2       & Argentina       &  39.934.100 \\ \midrule 
3       & Colômbia        &  46.741.100 \\ \midrule 
4       & Paraguai        &   9.694.200 \\ \midrule 
5       & Uruguai         &   3.350.500 \\ \midrule 
6      & Chile           &  16.803.000 \\ \bottomrule
\end{tabular}
}{%
  \fdireta{wikipedia:2011:america_sul}%
  \nota{Esta é uma nota, que diz que os dados são baseados na
  regressão linear.}%
  \nota[Anotações]{Uma anotação adicional, que pode ser seguida de várias
  outras, porém são opcionais.}%
  }
\end{table}

\begin{table}[htb]
\caption{Editores de Texto Livres}
\label{quadro:editores_texto_livres_novo}
\centering
\begin{tabular}{|l|l|r|}        \hline
Editor     & Multiplataforma & Específico para Latex \\ \hline
Kwriter    & Sim             & Não                   \\
Texmaker   & Sim             & Sim                   \\
Kile       & Sim             & Sim                   \\
Geany      & Sim             & Não                   \\ \hline
\end{tabular}
\end{table}

\begin{figure}[htb]
 \caption{Logomarca da USP}
 \label{fig:logomarca_usp_novo}
 \centering
 \includegraphics[scale=0.3]{usp-logo.pdf}
 \fdireta{usp:logo}
\end{figure}



Este documento explica brevemente como trabalhar com a classe \LaTeX~\textit{icmc} para confeccionar trabalhos acadêmicos seguindo as normas da \sigla{ABNT}{Associação Brasileira de Normas Técnicas} e as \aspas{\textit{Diretrizes para apresentação de dissertações e teses da USP: documento eletrônico e impresso. Parte I (ABNT)}}, publicado pelo \sigla{SIBi}{Sistema Integrado de Bibliotecas} USP. O presente manual também atende as exigências prevista no regimento do Programa de Pós-graduação em \sigla{CCMC}{Ciências da Computação e Matemática Computacional} do \sigla{ICMC}{Instituto de Ciências Matemáticas e de Computação} da \sigla{USP}{Universidade de São Paulo}.


A classe \textit{icmc} foi construída com base na última versão da classe \textit{abntex2} e do pacote \textit{abntex2cite}. Portanto, este documento exemplifica a elaboração de trabalho
acadêmico (tese, dissertação e outros do gênero) produzido conforme a ABNT NBR
14724:2011 \textit{Informação e documentação - Trabalhos acadêmicos - Apresentação}.

Assim, é altamente recomendável que seja consultada a documentação do \textit{abntex2}\footnote{https://code.google.com/p/abntex2/downloads/list}. A classe \textit{abntex2} foi desenvolvida para facilitar a escrita de documentos seguindo as normas da ABNT no ambiente \LaTeX\;\cite{frasson:2005:classe_abnt}.

Todo o trabalho de pesquisa e ajustes da presente classe \LaTeX~\emph{icmc} foram feitos pelo aluno mestrado do Programa de Pós-graduação em Ciência da Computação e Matemática Computacional, Humberto Lidio Antonelli, durante a confecção da sua monografia de qualificação.

O requisito básico para utilização da classe \textit{icmc} é criar um documento desta classe com o comando
\comando{documentclass[@parameters]\{icmc\}} e ter, no diretório de trabalho, o arquivo \emph{icmc.cls} presente. Entretanto, recomenda-se fortemente manter a estrutura de diretório inicial fornecida por este modelo. Além disso, para que o documento esteja em conformidade com as normas exigidas pelo programa de Pós-Graduação, o \textbf{projeto deve ser compilado utilizando \textit{XeLaTeX} ou \textit{LuaLaTeX}}. Esse processo de compilação é necessário para que as fontes externas utilizadas para gerar a capa sejam incluídas.

Os parâmetros possíveis utilizados pelo \comando{documentclass} são:
\begin{description}
\item[qualificacao] Exclusivamente para monografias de qualificação em geral;
\item[mestrado / doutorado] Identifica o curso ao qual o aluno pertence, sendo utilizado apenas uma das duas opcões disponíveis. O valor padrão é \textbf{doutorado};
\item[pre-defesa / pos-defesa] Identifica a situação do documento (exceto para qualificação), sedo necessário apenas uma das duas opções. O valor padrão é \textbf{pos-defesa};
\item[french, spanish, english, brazil] Adiciona o idioma para correta hifenização correta no documento. O último idioma declarado é o principal do documento. O valor padrão é \textbf{brazil}.
\end{description}